%% 导言区 —— SY02-43《一般线性群的局部 Langlands 对应》中译注本
%% 底本：Zijian Yao, "Notes on Local Langlands for GLn",
%%       Harvard 学习讨论班讲义（2017 年秋），36 页，pdfTeX 排版。
%% 编译：xelatex main.tex（三遍，待目录稳定）
\usepackage[UTF8, fontset=windows, heading=true]{ctex}
\usepackage{amsmath}
\usepackage{mathtools}
\usepackage{amssymb}
\usepackage{amsthm}
\usepackage{bm}
\usepackage[a4paper,left=2.6cm,right=2.6cm,top=2.6cm,bottom=2.6cm]{geometry}
\usepackage{graphicx}
\usepackage{xcolor}
\usepackage[most]{tcolorbox}
\usepackage{enumitem}
\usepackage{booktabs}
\usepackage{array}
\usepackage[font=small,labelfont=bf,labelsep=quad]{caption}
\usepackage[colorlinks=true,
            linkcolor=blue!45!black,
            citecolor=blue!45!black,
            urlcolor=blue!45!black]{hyperref}

\setlength{\emergencystretch}{3em}
\allowdisplaybreaks[3]
\linespread{1.12}
\renewcommand{\figurename}{图}
\renewcommand{\tablename}{表}
\ctexset{
  section/name       = {},
  section/number     = \arabic{section},
  subsection/name    = {},
  subsection/number  = \thesection.\arabic{subsection},
  subsubsection/name = {},
}

%% ---------- 定理族：底本中 定理/引理/推论/命题 共用一个按节重排的计数器 ----------
\theoremstyle{plain}
\newtheorem{nbthm}{定理}[section]
\newtheorem{theorem}[nbthm]{定理}
\newtheorem{lemma}[nbthm]{引理}
\newtheorem{corollary}[nbthm]{推论}
\newtheorem{proposition}[nbthm]{命题}
\newtheorem*{theorem*}{定理}

\theoremstyle{definition}
\newtheorem*{definition}{定义}
\newtheorem*{example}{例}
\newtheorem*{remark}{注}
\newtheorem*{corollarydefinition}{推论-定义}

\renewcommand{\proofname}{证明}

%% ---------- 「本章要点」（译者所加） ----------
\newtcolorbox{keytakeaways}[1][本章要点]{
  enhanced,
  colback  = blue!3!white,
  colframe = blue!35!black,
  boxrule  = 0.4pt,
  arc      = 2pt,
  left     = 7pt, right = 7pt, top = 5pt, bottom = 5pt,
  title    = {#1\,{\normalfont\scriptsize（译者所加）}},
  fonttitle= \bfseries\small,
  before skip = 10pt, after skip = 10pt,
}

%% ---------- 译注（译者所加） ----------
\newif\ifshowyizhu
\showyizhutrue
\newcounter{yizhu}
\makeatletter
\newenvironment{yizhu}[1][]{%
  \ifshowyizhu
    \refstepcounter{yizhu}%
    \par\addvspace{5pt}\noindent
    \begin{tcolorbox}[enhanced, breakable,
      colback  = orange!4!white,
      colframe = orange!55!black,
      boxrule  = 0.3pt,
      arc      = 1.5pt,
      left     = 6pt, right = 6pt, top = 4pt, bottom = 4pt,
      title    = {\small\bfseries 译注~\theyizhu
                  \if\relax\detokenize{#1}\relax\else\ ——#1\fi}]%
    \small
  \else
    \setbox\z@\vbox\bgroup\small
  \fi
}{%
  \ifshowyizhu
    \end{tcolorbox}%
  \else
    \egroup
  \fi
}
\makeatother

%% 底本用上标数字标脚注，译文以 \en{n} 复现并集中于文末「原注」
\newcommand{\en}[1]{\textsuperscript{#1}}
